% Open question on the stem model, written for the supervisors (standalone explanation).
% Source of the numbers: docs/TODO.md (open questions) and docs/notes/2026-09-30.md.
%
% Included by main.tex; uses only packages loaded there (amsmath, booktabs).


% ======================================= BEGIN CONTENT =======================================
\section{Open question: stiffness of the simulated stem}
\label{sec:stem-tradeoff}

\paragraph{Status (6 October 2026).} Answered: option (d) below, a chain of rigid segments in PhysX, was
built and compared with the cable (Section~\ref{sec:stem-comparison}); the supervisors chose it for the first
training stage. This section is kept as the record of why the cable was not used.

\subsection{Summary}

The simulated stem cannot have all three of the following at once: stretch and shear stiffness as high as
in a real stem, correct bending behaviour, and a simulation cheap enough for reinforcement learning. With
physically realistic stiffness values the solver does not resolve the bending within the available
computing budget, and the stem ends up about 23 times too soft and falls over. For now I make the stem
artificially softer against stretching and shearing, which restores correct bending at the cost of a stem
that can stretch slightly. The question for discussion is whether this workaround is acceptable, or whether
the stem should be modelled differently.

\subsection{How the stem is simulated}

The stem is modelled with the ``cable'' of the physics engine Newton, which is what Isaac Lab offers for
deformable linear objects. The cable is a chain of rigid segments (here 20 segments of 2\,cm for a stem of
0.4\,m). Neighbouring segments are connected by a joint with four springs, one for each way the stem can
deform:
\begin{itemize}
  \item \textbf{bending}, stiffness $EI/l$ per joint,
  \item \textbf{twisting} about the stem axis, stiffness $GJ/l$,
  \item \textbf{stretching} along the stem axis, stiffness $EA/l$,
  \item \textbf{shearing}, where neighbouring cross-sections slide sideways against each other, stiffness
        $GA/l$.
\end{itemize}
Here $l$ is the segment length, $A$ the cross-section area, $I$ and $J$ its area moments, and $E$, $G$ the
moduli. A real thin stem bends easily, but it practically does not stretch or shear. In the model this
means that the stretch and shear springs are far stiffer than the bending springs.

The solver (Vertex Block Descent, VBD) does not solve the equations of motion exactly. At every time step
it starts from a guess and improves it a fixed number of times. In each improvement pass it moves one
segment at a time, while the neighbours of that segment are held fixed. The computing effort per time step
is the number of substeps times the number of passes per substep; I call this product the \emph{cost}.

\subsection{The problem}

In the exact solution the two kinds of springs do not conflict. The stiff springs decide what the stem
cannot do (stretch, shear), and the soft bending springs decide what it does within those rules. The
difficulty is reaching that solution.

The error that the solver reduces can be pictured as a long, narrow valley. Its walls are steep, because
stretching or shearing the stem costs a lot of energy. Its floor is almost flat, because bending costs
little. The solver drops down the steep walls within a few passes and then creeps along the flat floor.
When it is stopped after a fixed number of passes, stretch and shear are satisfied, but the bending is
not.

The mechanical reason is the one-segment-at-a-time update. To bend the stem at one joint, everything beyond
that joint has to swing together. A single segment that turns on its own moves its two ends sideways
against its fixed neighbours, which the stiff shear springs resist immediately. Each segment may therefore
turn only by a tiny amount per pass, and a bend is passed along the chain in small steps.

The resulting error is not random noise. It appears as a wrong bending stiffness, and it can go in either
direction: the stem is too soft where it should hold itself up (it has not caught up with its own restoring
force), and too stiff where it should yield to a push (it has not moved far enough yet).

\subsection{Why it cannot be tuned away}

What matters is not the size of the stiffness values but the ratio between shear (or stretch) stiffness
and bending stiffness. For a round stem this ratio is roughly
\begin{equation*}
  \frac{\text{shear stiffness}}{\text{bending stiffness}} \;\propto\;
  \frac{G}{E}\left(\frac{l}{d}\right)^{2},
\end{equation*}
with segment length $l$ and stem diameter $d$. This is my own estimate; it agrees with the measurements
below, but it was tested on one geometry only. Two consequences follow. First, scaling all moduli by the
same factor, or changing units, changes nothing. Second, for a physical material $G/E$ is fixed, so the
only remaining parameter is the segment length relative to the diameter. With physical moduli the segments
would have to be about 1\,mm long, which means several hundred segments per stem.

Solving bending and shear in two separate steps does not help either. They are not separate unknowns: the
unknowns are the position and orientation of every segment, and both springs depend on them. A bending
step followed by a shear step would partly undo each other, and the alternation converges as slowly as
before.

\subsection{Measurements}

All measurements use the stem alone, clamped at its base, with a time step of 10\,ms. The stem has
placeholder values: length 0.4\,m, diameter 8\,mm, 20 segments, density 1000\,kg/m$^3$, bend modulus
$5\cdot10^{8}$\,Pa. Stretch and shear moduli are given as multiples of the bend modulus; a value of 1 is
the physical order of magnitude. Three tests are used, each compared with a value computed by hand for the
same chain of segments:
\begin{itemize}
  \item \textbf{Sag:} the stem is clamped horizontally and sags under its own weight. Reported is the tip
        sag divided by the hand value, so 1 is correct. (For many segments the hand value approaches the
        beam formula $qL^{4}/(8EI)$.)
  \item \textbf{Kick:} the stem stands upright and its tip is kicked sideways with 1\,m/s. It should swing
        and return to upright.
  \item \textbf{Push:} the stem stands upright and its last segment is pushed sideways with a constant
        force of 0.05\,N. The expected deflection is about 10.2\,mm: 9.85\,mm by hand without gravity
        (approaching $FL^{3}/(3EI)$ for many segments), plus about 4\,\% because the stem's own weight
        makes it lean further.
\end{itemize}

\begin{table}[h]
  \centering
  \caption{Sag divided by the hand value when stretch and shear are lowered together (1 is correct).}
  \label{tab:stem-together}
  \begin{tabular}{lccc}
    \toprule
    Stretch = shear modulus & Cost 80 & Cost 160 & Cost 800 \\
    \midrule
    1 (physical)  & not run      & 23 (stem falls over when kicked) & 15, still moving \\
    0.1           & not run      & 22   & 7     \\
    0.01          & not run      & 3.7  & 1.28  \\
    0.001         & 1.30 to 1.44 & 1.06 & 1.003 \\
    \bottomrule
  \end{tabular}
\end{table}

\begin{table}[h]
  \centering
  \caption{Effect of the stretch modulus alone. Sag ratio = sag divided by the hand value. Shear modulus 0.001, cost 80 (8 substeps, 10 passes).
           The push deflection should be about 10.2\,mm.}
  \label{tab:stem-stretch}
  \begin{tabular}{lcllc}
    \toprule
    Stretch modulus & Sag ratio & Kick & Push deflection & Stretch per N \\
    \midrule
    1               & 1.05    & blows up          & not run                     & 0.016\,mm \\
    0.1             & 1.012   & returns, 5.05\,Hz & 8.56\,mm (16\,\% too stiff) & 0.16\,mm  \\
    0.03            & 1.009   & returns, 4.75\,Hz & 9.91\,mm (3\,\% too stiff)  & 0.5\,mm   \\
    0.01 (in use)   & 1.008   & returns, 4.67\,Hz & 10.27\,mm (1\,\% too soft)  & 1.6\,mm   \\
    0.001           & not run & not run           & 10.38\,mm (2\,\% too soft)  & 16\,mm    \\
    \bottomrule
  \end{tabular}
\end{table}

\begin{table}[h]
  \centering
  \caption{More solver work instead of softer stretch (stretch 0.1, shear 0.001). Push deflection,
           expected about 10.2\,mm.}
  \label{tab:stem-solver-work}
  \begin{tabular}{lcc}
    \toprule
    Substeps $\times$ passes & Cost & Push deflection \\
    \midrule
    $8 \times 10$  & 80  & 8.56\,mm \\
    $16 \times 5$  & 80  & 9.40\,mm \\
    $8 \times 15$  & 120 & 9.11\,mm \\
    $12 \times 10$ & 120 & 9.49\,mm \\
    $8 \times 20$  & 160 & 9.39\,mm \\
    $16 \times 10$ & 160 & 9.87\,mm \\
    $8 \times 40$  & 320 & 9.86\,mm \\
    \bottomrule
  \end{tabular}
\end{table}

The tables show three things.
\begin{itemize}
  \item With physical values the stem is far too soft, and more solver work alone does not fix it within
        any affordable budget (Table~\ref{tab:stem-together}).
  \item The main cause is the shear spring. Lowering only the shear modulus to 0.001 already gives the
        correct sag (first row of Table~\ref{tab:stem-stretch}). The stretch spring must be lowered as well: at full stretch
        stiffness the kick test becomes unstable.
  \item Which test is run decides what is seen. With a stretch modulus of 0.1 the sag test looks correct,
        but the push on the upright stem is 16\,\% too stiff. In the upright stem gravity acts along the
        stem and loads the stiff stretch springs; in the horizontal sag test it does not. Without gravity
        the same setting gives 10.06\,mm in the push test.
\end{itemize}

One further observation concerns damping. Damping the stretch and shear springs slows the solver down: with
all four springs damped, the kick test is fine at cost 160, but the stem no longer returns to upright at
cost 80 and the simulation becomes unstable at cost 40. With damping on bending and twisting only, the
test passes at cost 80. For reference, the computing time with 1024 stems in parallel is 8.6\,ms per
10\,ms time step at cost 80 and 16.8\,ms at cost 160.

\subsection{Current workaround and what it costs}

The setting in use is: stretch modulus 0.01 and shear modulus 0.001 times the bend modulus, damping on
bending and twisting only, 8 substeps with 10 passes each (cost 80). With this setting the sag and the push
are within 1\,\% of the expected values, and after a kick the stem swings at 4.67\,Hz with a damping ratio
of 0.06 and returns to upright.

The price is the following.
\begin{itemize}
  \item \textbf{The stem can stretch.} It elongates by 1.6\,mm per newton of pull along its axis. A real
        stem practically does not stretch.
  \item \textbf{The stem is about 10\,\% softer in bending} than its bending stiffness alone implies,
        because the soft shear spring adds deflection of its own (8\,\% for a force at the tip). This
        could be compensated by a 10\,\% higher bend modulus.
  \item \textbf{Stretch and shear vibrations are not damped} by the material model, only by the small
        numerical damping of the solver (damping ratio 0.01 to 0.02).
\end{itemize}

The argument for accepting this is that a thin stem hardly stretches or shears in reality, so the exact
stiffness in these directions has little influence on how it moves, as long as it is high enough. The task
(moving a point on the stem to a target) and the damage limit (maximum curvature) both depend on bending.

\subsection{Limits of these results}

The measurements were made with one stem geometry and placeholder material values, without contact and
without the robot. An accuracy statement holds only for the load cases that were tested (sag, kick, push);
the push test revealed an error that the other two had not shown. The usable stiffness ratios depend on
segment length and diameter, so they must be checked again when the stem parameters change, when they are
randomized for training, and in the coupled simulation with the robot.

\subsection{Options}

\begin{table}[h]
  \centering
  \caption{Options. Bending error is given for the sag and the push test.}
  \label{tab:stem-options}
  \begin{tabular}{llclll}
    \toprule
     & Change & Cost & Bending error & Stretch per N & Price \\
    \midrule
    (a) & stretch 0.01 (in use) & 80  & 1\,\% / 1\,\%  & 1.6\,mm  & stretchy stem \\
    (b) & stretch 0.03          & 80  & 1\,\% / 3\,\%  & 0.5\,mm  & compromise \\
    (c) & stretch 0.1, $16 \times 10$ passes & 160 & 1\,\% / 3\,\% & 0.16\,mm & half the speed \\
    (d) & different stem model          & unknown & none by construction & none & new model \\
    \bottomrule
  \end{tabular}
\end{table}

Option (d) replaces the cable by a chain of rigid segments connected by joints that can only bend and
twist, each with a rotational spring of stiffness $EI/l$. Such a chain is described by its joint angles. It
cannot stretch or shear by construction, which is exactly the assumption of thin-rod theory, so the stiff
springs that cause the problem do not exist. This model would run in the same solver as the robot arm, so
the coupling of two solvers would no longer be needed, and clamping the base and joint damping are
standard features there. The price is that the model has to be built and validated.

More generally, the difficulty comes from the combination of this stem description (every segment is free
and held to its neighbours by stiff springs) with this kind of solver (one segment at a time). Changing
either one removes it: describing the stem by joint angles, as in (d), or solving all segments of the chain
at once with a direct solver for rods, which Newton's VBD solver does not offer. Softening the springs only
makes the problem mild enough for the solver at hand.

\subsection{Questions for discussion}

\begin{enumerate}
  \item Is softening the stretch and shear stiffness an accepted practice for this kind of rod model, and
        how soft is acceptable for a plant stem?
  \item Is there a solver setting or formulation that handles the stiff directions? (The solver's
        ``compliant'' constraint mode made no difference in my tests.)
  \item Is the rigid-segment chain (d) the better model for a stiff, thin stem? Domain randomization of
        stem stiffness and diameter will change the stiffness ratio from plant to plant, which makes a
        single softened setting harder to justify.
\end{enumerate}
% ======================================== END CONTENT ========================================
